%%%%%%%%%%%%%%%%%%%%%%%%%%%%%%%%%%%%%%%%%%%%%%%%%%%%%%%%%%%%%%%%%%%%%%
% Lớp 10 | Đại số tổ hợp | Hoán vị. Chỉnh hợp. Tổ hợp
% Dạng 0D8?2-4: Bài toán chọn người
% Số câu: 46  (sinh tự động bởi tach_cau.py; từ nay thầy sửa trực tiếp file này)
%%%%%%%%%%%%%%%%%%%%%%%%%%%%%%%%%%%%%%%%%%%%%%%%%%%%%%%%%%%%%%%%%%%%%%

%%%%%%%%%%%%%%%% TRẮC NGHIỆM 4 PHƯƠNG ÁN %%%%%%%%%%%%%%%%

% nguồn: MaHoaNH_TracNghiem_Lop10
\begin{ex}%[0D8N2-4]%[ma:fd43a0d0]
	Tổ một có $11$ học sinh. Cô giáo muốn chọn $6$ bạn trong tổ một đi trực nhật khu tự quản. Hỏi cô giáo có bao nhiêu  cách chọn?
	\choice 
		 {$\mathrm{A}_{11}^{6}$}
		{\True $\mathrm{C}_{11}^{6}$} 
		 {$16$}
		{$10$} 
	 \loigiai{
		Mỗi cách chọn $6$ bạn trong tổ một đi trực nhật khu  tự quản là một tổ hợp chập $6$ trong $11$ phần tử, do đó số cách chọn là $\mathrm{C}_{11}^{6}$.
	 }
\end{ex}

% nguồn: MaHoaNH_TracNghiem_Lop10
\begin{ex}%[0D8N2-4]%[ma:fa0e1c79]
	Tên $15$ học sinh được ghi vào $15$ tờ giấy để vào trong hộp. Chọn tên $4$ học sinh để cho đi du lịch. Hỏi có bao nhiêu cách chọn các  học sinh?
	\choice 
		{$4!$} 
		 {$15!$}
		 {\True $1365$}
		 {$32760$}
	\loigiai{ 
		Số cách chọn $4$ học sinh để cho đi du lịch là số cách  chọn $4$ trong $15$ tờ giấy trong hộp, do đó có $\mathrm{C}_{15}^{4} = 1365$ cách.
	 }
\end{ex}

% nguồn: MaHoaNH_TracNghiem_Lop10
\begin{ex}%[0D8N2-4]%[ma:38bdfdfa]
	Một ban cán sự lớp có $3$ người  gồm $1$ lớp trưởng, $1$ bí thư và $1$ ủy viên. Tính số cách tạo một ban cán sự lớp từ một nhóm $10$ học sinh.
	 \choice
		 {$120$}
		 {\True $720$}
		{$1000$} 
		{$30$} 
	 \loigiai{
		Số cách tạo  một ban cán sự lớp từ một nhóm $10$ học sinh là $\mathrm A_{10}^3=720$.
	} 
\end{ex}

% nguồn: MaHoaNH_TracNghiem_Lop10
\begin{ex}%[0D8N2-4]%[ma:9b908ecf]
	Lớp $10A$ có $20$ học sinh. Hỏi giáo viên chủ nhiệm lớp $10A$ có tất cả  bao nhiêu cách chọn ra $5$ học sinh tham gia cuộc thi văn nghệ của nhà trường?
	\choice 
		{\True $15504$} 
		 {$120$}
		 {$3003$}
		{$1860480$} 
	 \loigiai{
		Số cách chọn là $\mathrm{C}_{20}^{5}=15504$. 
	 }
\end{ex}

% nguồn: MaHoaNH_TracNghiem_Lop10
\begin{ex}%[0D8N2-4]%[ma:faaef9c9]
	Một tổ có $6$ học sinh nam và $5$ học sinh nữ. Có bao nhiêu cách chọn một học sinh nam và một  học sinh nữ để đi tập văn nghệ?
	 \choice
		{\True  $30$}
		 {$\mathrm{C}_{11}^2$}
		{$11$} 
		{$\mathrm{A}_{11}^2$} 
	\loigiai{ 
		Có $6$ cách chọn $1$ học  sinh nam từ $6$ học sinh nam.\\
		Ứng với mỗi cách chọn $1$ học  sinh nam có $5$ cách chọn $1$ học sinh nữ từ $5$ học sinh nữ.\\
		Theo quy tắc nhân có $6\cdot 5=30$ cách chọn một học sinh nam và một học sinh nữ để đi  tập văn nghệ.
	} 
\end{ex}

% nguồn: MaHoaNH_TracNghiem_Lop10
\begin{ex}%[0D8N2-4]%[ma:6f46ced7]
	Lớp 10K có $ 32 $  bạn nam và $ 8 $ bạn nữ. Có bao nhiêu cách chọn một bạn để làm lớp trưởng?
	\choice 
		 {$ 8 $}
		{$ 256 $} 
		{$  32 $}
		{\True  $ 40 $}
	\loigiai{ 
		Để chọn $ 1 $ bạn làm lớp trưởng ta có hai  phương án sau
		\begin{itemize}
			\item Phương án 1: Chọn $  1 $ bạn nam có $ 32 $ cách.
			 \item Phương án 2: Chọn $ 1 $ bạn nữ có $ 8 $ cách.
		\end{itemize}
		 Theo quy tắc cộng ta có $ 32+8=40 $ cách chọn $ 1 $ bạn làm lớp trưởng.
	} 
\end{ex}

% nguồn: MaHoaNH_TracNghiem_Lop10
\begin{ex}%[0D8H2-4]%[ma:6c0d78f9]
	Một lớp học gồm có $40$ học sinh, cần chọn một ban cán sự lớp gồm $3$ người với  các chức vụ Lớp trưởng, lớp phó học tập và lớp phó lao động. Hỏi có bao nhiêu cách chọn?
	\choice 
		 {$3!$}
		 {\True $\mathrm{A}_{40}^3$}
		{$\mathrm{C}_{40}^3$} 
		{$40!$} 
	\loigiai{ 
		Vì chọn ra $3$ người từ $40$ người sau đó  sắp xếp chức vụ cho $3$ người đó nên số cách chọn là $\mathrm{A}_{40}^3$.
	 }
\end{ex}

% nguồn: MaHoaNH_TracNghiem_Lop10
\begin{ex}%[0D8H2-4]%[ma:4817632b]
	Một giải thể thao chỉ có ba giải là nhất, nhì, ba.  Trong số $20$ vận động viên đi thi, có bao nhiêu khả năng mà ba người có thể được ban tổ chức trao giải nhất, nhì, ba?
	\choice 
		 {$1\,140$}
		{\True  $6\,840$}
		{$1$} 
		{$3$} 
	\loigiai{ 
		Số  khả năng là $\mathrm{A}_{20}^3=6\,840$.
	 }
\end{ex}

% nguồn: MaHoaNH_TracNghiem_Lop10
\begin{ex}%[0D8H2-4]%[ma:97311088]
	Trong một dạ hội cuối năm ở một cơ quan, ban tổ chức phát ra $100$ vé xổ số đánh số từ $1$ đến $100$ cho $100$ người. Xổ số có $4$ giải: $1$ giải  nhất, $1$ giải nhì, $1$ giải ba, $1$ giải tư. Kết quả là việc công bố ai trúng giải nhất, giải nhì, giải ba, giải tư. Hỏi có bao nhiêu kết quả có thể nếu biết rằng người giữ vé số $47$ được giải nhất?
	\choice 
		 {\True $944\,109$}
		 {$941\,409$}
		{$941\,094$} 
		{$156\,849$} 
	 \loigiai{
		Vì người giữ vé số $47$ trúng giải nhất nên mỗi  kết quả ứng với một chỉnh hợp chập $3$ của $99$ phần tử, do đó ta có $\mathrm{A}_{99}^3=941\,094$ kết quả.
	} 
\end{ex}

% nguồn: MaHoaNH_TracNghiem_Lop10
\begin{ex}%[0D8H2-4]%[ma:ef293a9e]
	Trong một buổi  khiêu vũ có $20$ nam và $18$ nữ. Hỏi có bao nhiêu cách chọn ra một đôi nam nữ để khiêu vũ?
	 \choice
		 {$\mathrm{C}_{38}^{2}$}
		 {$\mathrm{A}_{38}^{2}$}
		{$\mathrm{C}_{20}^{2}\cdot \mathrm{C}_{18}^{1}$} 
		{\True  $\mathrm{C}_{20}^{1}\cdot \mathrm{C}_{18}^{1}$}
	\loigiai{ 
		Có $\mathrm{C}_{20}^{1}\cdot \mathrm{C}_{18}^{1}$ cách chọn ra một đôi nam nữ để khiêu vũ. 
	 }
\end{ex}

% nguồn: MaHoaNH_TracNghiem_Lop10
\begin{ex}%[0D8H2-4]%[ma:b61ff948]
	Hội đồng quản trị của công ty X gồm $10$ người. Hỏi có bao nhiêu cách bầu  ra ba người vào ba vị trí chủ tịch, phó chủ tịch và thư kí, biết khả năng mỗi người là như nhau?
	\choice 
		 {$728$}
		{$723$} 
		{\True $720$} 
		 {$722$}
	 \loigiai{
		Số cách bầu ra ba người vào ba vị trí chủ tịch, phó chủ tịch và thư kí là số cách chọn $3$ người trong $10$  người và chọn vị trí cho $3$ người được chọn. Do đó, số cách bầu là $\mathrm{A}_{10}^{3} = 720$.
	} 
\end{ex}

% nguồn: MaHoaNH_TracNghiem_Lop10
\begin{ex}%[0D8H2-4]%[ma:d4c035bb]
	Một tổ gồm $ 8 $ nam và $ 4 $ nữ. Hỏi có bao nhiêu cách chọn  $ 4 $ em đi trực sao cho có đúng $ 1 $ nữ?
	\choice 
		{\True $  \mathrm{C}_8^3\cdot \mathrm{C}_4^1 $}
		{\True $ \mathrm{C}_{12}^4  $}
		{$ \mathrm{C}_4^1 $} 
		{$ \mathrm{C}_8^3+\mathrm{C}_4^1  $}
	 \loigiai{
		 Chọn $ 1 $ học sinh nữ trong $ 4 $ học sinh nữ và $ 3 $ học sinh nam trong $ 8 $ học sinh nam có $ \mathrm{C}_8^3\cdot \mathrm{C}_4^1 $ cách.
	} 
\end{ex}

% nguồn: MaHoaNH_TracNghiem_Lop10
\begin{ex}%[0D8H2-4]%[ma:d517ef81]
	Trong một trận chung kết bóng đá cần phải đá luân lưu  $11$ mét để phân định thắng thua. Huấn luyện viên cần trình với trọng tài một danh sách $3$ cầu thủ trong $7$ cầu thủ đang có trên sân để lần lượt theo thứ tự đá đủ $3$ quả sút luân lưu (mỗi cầu thủ đá đúng một lần). Huấn luyện viên có tất cả bao nhiêu cách chọn?
	 \choice
		{$70$} 
		 {$2\,187$}
		 {$823\,543$}
		{\True $210$} 
	 \loigiai{
		Mỗi cách chọn ra $3$ cầu thủ trong $7$ cầu thủ để thực hiện đá luân lưu (theo thứ tự) là một chỉnh hợp chập $3$ của $7$ phần tử.  \\
		Vậy số cách chọn  là $\mathrm{A}_7^3=210$ cách.
	 }
\end{ex}

% nguồn: MaHoaNH_TracNghiem_Lop10
\begin{ex}%[0D8H2-4]%[ma:4ee4dc33]
	Một lớp học có $40$ học sinh gồm $15$ nam và  $25$ nữ. Giáo viên cần chọn $3$ học sinh tham gia lao động. Hỏi có bao nhiêu cách chọn khác nhau?
	\choice 
		 {\True $9\,880$}
		{$59\,280$} 
		 {$2\,300$}
		 {$455$}
	 \loigiai{
		Mỗi  cách chọn ra $3$ học sinh để tham gia lao động từ $40$ học sinh là một tổ hợp chập $3$ của $40$ phần tử.\\
		Vậy  có $\mathrm{C}_{40}^3=9\,880$ cách.
	 }
\end{ex}

% nguồn: MaHoaNH_TracNghiem_Lop10
\begin{ex}%[0D8H2-4]%[ma:ed4aa2b2]
	Một tổ có $6$ học sinh nam và $9$  học sinh nữ. Hỏi có bao nhiêu cách chọn $1$ học sinh nam và $1$ học sinh nữ đi lao động?
	 \choice
		 {$\mathrm{C}_6^1+\mathrm{C}_9^1$}
		 {$\mathrm{C}_6^1\cdot\mathrm{C}_{15}^1$}
		 {$\mathrm{C}_6^1+\mathrm{C}_{15}^1$}
		{\True  $\mathrm{C}_6^1\cdot \mathrm{C}_9^1$}
	\loigiai{ 
		Chọn $1$ học sinh nam từ $6$ học sinh nam có  $\mathrm{C}_6^1$ cách chọn.\\
		Chọn $1$ học sinh nữ từ $9$ học sinh nữ có $\mathrm{C}_9^1$ cách chọn.\\ 
		Vậy có $\mathrm{C}_6^1\mathrm{C}_9^1$ cách chọn $2$ học sinh đi lao động trong đó có đúng $1$ nam và $1$  nữ.
	} 
\end{ex}

% nguồn: MaHoaNH_TracNghiem_Lop10
\begin{ex}%[0D8H2-4]%[ma:2dc7c8be]
	Tổ $1$ lớp 10A có $3$ học sinh nữ và $5$ học sinh nam. Có bao nhiêu  cách xếp các học sinh tổ $1$ thành một hàng dọc sao cho không có hai bạn nữ nào đứng cạnh nhau?
	 \choice
		 {$4\,800$}
		 {$3\,600$}
		 {\True $7\,200$}
		 {$9\,600$}
	 \loigiai{
		Xếp các  bạn nam trước có $5!$ cách, tạo ra $6$ khoảng trống.\\
		Xếp $3$ bạn  nữ vào các khoảng trống đó có $\mathrm{A}_5^3$ cách.\\
		Do đó, có $5!\cdot \mathrm{A}_5^3=7\,200$ cách sắp xếp thỏa  mãn yêu cầu.
	} 
\end{ex}

% nguồn: MaHoaNH_TracNghiem_Lop10
\begin{ex}%[0D8H2-4]%[ma:63ffb9db]
	Một lớp có $25$ học sinh nam và $20$ học sinh  nữ, số cách chọn $3$ em học sinh đi dự đại hội đoàn trong đó có nhiều nhất $1$ em nữ là
	 \choice
		{$12\,102$} 
		{\True $8\,300$} 
		 {$6\,000$}
		 {$1\,200$}
	\loigiai{ 
		Chọn $3$ em có $1$ em nữ có  $\mathrm{C}_{25}^2\cdot\mathrm{C}_{20}^1$ cách.\\
		Chọn $3$ em trong đó không có em nữ nào  có $\mathrm{C}_{25}^3$ cách.\\
		Số cách chọn $3$  em học sinh trong đó có nhiều nhất $1$ em nữ là\\ $\mathrm{C}_{25}^2\cdot\mathrm{C}_{20}^1+\mathrm{C}_{25}^3=8\,300$ cách.
	} 
\end{ex}

% nguồn: MaHoaNH_TracNghiem_Lop10
\begin{ex}%[0D8H2-4]%[ma:c3aafc50]
	Có $10$ y tá và $3$ bác sĩ. Hỏi có bao nhiêu cách lập nhóm công tác gồm $1$ bác sĩ  làm trưởng đoàn, $1$ y tá làm phó đoàn và $5$ y tá làm thành viên?
	\choice 
		 {$8\,730$}
		{\True $3\,780$} 
		 {$3\,870$}
		{$7\,830$} 
	 \loigiai{
		Chọn  một bác sĩ làm trưởng đoàn có $3$ cách chọn.\\
		 Chọn một y tá làm phó đoàn có $10$ cách chọn.\\
		Chọn $5$ y tá làm thành  viên có $\mathrm{C}_9^5$ cách chọn.\\
		V  ậy có $3\cdot 10\cdot \mathrm{C}_9^5=3\,780$ cách chọn.
	 }
\end{ex}

% nguồn: MaHoaNH_TracNghiem_Lop10
\begin{ex}%[0D8H2-4]%[ma:8c1ff953]
	Lớp 12A5 của trường THPT Nguyễn Công Trứ có $15$ học sinh giỏi theo từng tổ như sau tổ $1$ có $5$ học sinh, tổ $2$ có $5$ học sinh và tổ $3$ có $5$ học sinh. Nhà trường cần chọn một đội tuyển gồm $10$ học sinh tham gia thi HSG cấp tỉnh. Tính số cách lập đội tuyển sao cho có học sinh cả ba tổ  và có nhiều nhất $2$ học sinh tổ $1$.
	\choice 
		{$59$} 
		{$50$} 
		 {$506$}
		{\True $500$} 
	 \loigiai{
		Ta có các trường hợp  sau\\
		Trường hợp  1: Có đúng $1$ học sinh tổ $1$, $9$ học sinh còn lại chọn từ tổ $1$ và tổ $2$ thì có $\mathrm{C}_5^1\cdot \mathrm{C}_{10}^9$ cách.\\
		Trường hợp 2: Có $2$ học sinh  tổ $1$, $8$ học sinh còn lại chọn từ tổ $1$ và tổ $2$ thì có $\mathrm{C}_5^2\cdot \mathrm{C}_{10}^8$ cách.\\
		 Vậy có $500$ cách.
	} 
\end{ex}

% nguồn: MaHoaNH_TracNghiem_Lop10
\begin{ex}%[0D8H2-4]%[ma:c415c1c8]
	Một công ty có $9$ nam và $13$ nữ trong  đó có đúng $2$ cặp vợ chồng. Công ty muốn chọn $1$ đội chơi tham dự teambuilding gồm $5$ người. Hỏi có bao nhiêu cách chọn ra $5$ người trong số $22$ người đó nhưng không có cặp vợ chồng nào?
	\choice 
		 {$24\,054$}
		{$24\,075$} 
		 {$24\,090$}
		 {\True $24\,072$}
	\loigiai{ 
		Ta có các trường hợp sau\\ 
		Trường hợp 1: Chọn $5$ người từ $18$ người không có cặp vợ chồng nào, trường hợp  này có $\mathrm{C}_{18}^5$ cách.\\
		Trường hợp 2:  chọn $1$ người từ $2$ cặp vợ chồng và $4$ người từ $18$ người thì có $\mathrm{C}_4^1\cdot \mathrm{C}_{18}^4$ cách.\\
		Trường hợp 3: chọn $2$ người từ $2$ cặp vợ chồng sao  cho không phải là một cặp và $3$ người từ $18$ người thì có $(\mathrm{C}_4^2-2)\mathrm{C}_{18}^3$ cách.\\
		Vậy có $24\,072$ cách. 
	} 
\end{ex}

% nguồn: MaHoaNH_TracNghiem_Lop10
\begin{ex}%[0D8H2-4]%[ma:9f03a496]
	Một nhóm $9$ người gồm ba đàn ông, bốn phụ nữ và hai đứa trẻ đi xem  phim. Hỏi có bao nhiêu cách xếp họ ngồi trên một hàng ghế sao cho mỗi đứa trẻ ngồi giữa hai phụ nữ và không có hai người đàn ông nào ngồi cạnh nhau?
	\choice 
		{$288$} 
		{\True $864$} 
		 {$24$}
		{$576$} 
	 \loigiai{
		Dùng phương pháp vách ngăn.\\ 
		 Xếp $4$ người phụ nữ có $4!$ cách.\\
		Khi đó có $3$ chỗ để xếp $2$ trẻ  em ngồi giữa $2$ người phụ nữ thì có $\mathrm{A}_3^2$ cách.\\
		Sau khi xếp $4$ người phụ nữ và $2$ trẻ em thì có $3$ chỗ để xếp $3$ người đàn ông sao cho không có $2$ người nào  ngồi cạnh nhau thi có $3!$.\\
		Vậy  có $3!\cdot \mathrm{A}_3^2\cdot 4!=864$ cách.
	 }
\end{ex}

% nguồn: MaHoaNH_TracNghiem_Lop10
\begin{ex}%[0D8H2-4]%[ma:f3aefa8c]
	Một đội văn nghệ của nhà trường gồm $4$ học sinh lớp $12A$, $3$ học sinh lớp $12B$  và $2$ học sinh lớp $12C$. Chọn ngẫu nhiên $5$ học sinh từ đội văn nghệ để biểu diễn trong lễ bế giảng năm học. Hỏi có bao nhiêu cách chọn sao cho lớp nào cũng có học sinh được chọn?
	\choice 
		{\True $98$} 
		 {$66$}
		 {$360$}
		 {$32$}
	\loigiai{ 
		Chọn ngẫu nhiên $5$ học sinh từ đội văn nghệ của  trường có $\mathrm{C}_9^5=126$ cách.\\
		Ta tính số cách chọn  $5$ học sinh sao cho không có đủ học sinh của ba lớp.\\
		{\bfseries Trường  hợp 1.} Chọn $5$ học sinh từ hai lớp $A$, $B$ có $\mathrm{C}_7^5=21$ cách.\\
		{\bfseries Trường hợp  2.} Chọn $5$ học sinh từ hai lớp $A$, $C$ có $\mathrm{C}_6^5=6$ cách.\\
		{\bfseries Trường  hợp 3.} Chọn $5$ học sinh từ hai lớp $B$, $C$ có $\mathrm{C}_5^5=1$ cách.\\
		Do đó có $21+6+1=28$ cách  chọn $5$ học sinh sao cho không có đủ học sinh của ba lớp.\\
		Vậy có $126-28=98$ cách chọn $5$ học sinh sao cho có đủ học sinh của ba  lớp.
	 }
\end{ex}

% nguồn: MaHoaNH_TracNghiem_Lop10
\begin{ex}%[0D8H2-4]%[ma:db4625f0]
	Có $14$ người gồm $8$ nam và $6$ nữ.  Có bao nhiêu cách chọn một tổ $6$ người trong đó có nhiều nhất $2$ nữ?
	 \choice
		 {$1\ 524$}
		{$472$} 
		{\True $1\  414$}
		{$3\ 003$} 
	\loigiai{ 
		Ta có các trường hợp  sau:
		\begin{itemize}
			\item Chọn  $6$ nam và không có nữ có $\mathrm{C}_8^6=28$ (cách).
			\item Chọn $1$ nữ và $5$ nam có  $\mathrm{C}_6^1\mathrm{C}_8^5=336$ (cách).
			\item Chọn  $2$ nữ $4$ nam có $\mathrm{C}_6^2\mathrm{C}_8^4=1\ 050$ (cách).
		\end{itemize}
		Theo quy tắc cộng ta có $28+336+1\ 050=1\ 414$ cách để chọn một  tổ có $6$ người trong đó có nhiều nhất $2$ nữ.
	 }
\end{ex}

% nguồn: MaHoaNH_TracNghiem_Lop10
\begin{ex}%[0D8V2-4]%[ma:3bac825c]
	Có hai  học sinh lớp $A$, ba học sinh lớp $B$ và bốn học sinh lớp $C$ xếp thành một hàng ngang sao cho giữa hai học sinh lớp $A$ không có học sinh nào lớp $B$. Hỏi có bao nhiêu cách xếp hàng như vậy?
	\choice 
		 {$80\,640$}
		 {$108\,864$}
		{\True $145\,152$} 
		 {$217\,728$}
	\loigiai{ 
		Xét các  trường hợp sau\\
		Trường hợp  1: Hai học sinh lớp A đứng cạnh nhau có $2!.8!$ cách.\\
		Trường hợp  2: Giữa hai học sinh lớp A có một học sinh lớp C có $2!\cdot \mathrm{A}_4^1\cdot 7!$ cách.\\
		Trường hợp 3: Giữa hai học sinh lớp A có hai học  sinh lớp C có $2!\cdot \mathrm{A}_4^2\cdot 6!$ cách.\\
		Trường hợp 4: Giữa hai học sinh lớp A có ba học sinh lớp  C có $2!\cdot \mathrm{A}_4^3\cdot 5!$ cách.\\
		Trường hợp 5: Giữa hai học sinh  lớp A có bốn học sinh lớp C có $2!\cdot \mathrm{A}_4^4\cdot 4!$ cách.\\
		Vậy  theo quy tắc cộng có $2!\left(8!+\mathrm{A}_4^17!+\mathrm{A}_4^26!+\mathrm{A}_4^35!+\mathrm{A}_4^44!\right)=145\,152$ cách.
	} 
\end{ex}


%%%%%%%%%%%%%%%% ĐÚNG/SAI %%%%%%%%%%%%%%%%

% nguồn: MaHoaNH_DungSai_Lop10
\begin{ex}%[0D8N2-4]%[ma:4dd2f413]
	Chọn khẳng định đúng trong các  khẳng định sau
	\choiceTF 
		{Có $24360$ cách chọn $3$ người đi  công tác từ một tổ có $30$ người}
		{Có $870$ cách chọn $2$ người đi công  tác từ một tổ có $30$ người}
		{\True Có $435$ cách chọn $2$ người đi công  tác từ một tổ có $30$ người}
		{\True Có $4060$ cách chọn $3$ người đi công tác từ một tổ có  $30$ người}
	\loigiai{ 
		\begin{itemize}
			\item Số cách chọn $3$ người bất kì trong $30$ là  $\mathrm{C}_{30}^3=4060$.
			 \item Số cách chọn $2$ người bất kì trong $30$ là $\mathrm{C}_{30}^2=435$.
		\end{itemize}
	} 
\end{ex}

% nguồn: MaHoaNH_DungSai_Lop10
\begin{ex}%[0D8H2-4]%[ma:69135128]
	Trong một buổi khiêu vũ có $20$ nam và $18$ nữ. Chọn khẳng  định đúng trong các khẳng định sau.
	\choiceTF 
		{\True Có $20$ cách  chọn ra một nam}
		{\True Có $18$ cách  chọn ra một nữ}
		{Có $703$ cách chọn ra một đôi nam nữ để  khiêu vũ}
		{\True Có $360$ cách chọn ra một  đôi nam nữ để khiêu vũ}
	\loigiai{ 
		Chọn một nam trong $20$ nam  có $\mathrm{C}_{20}^1=20$ cách.\\
		Chọn một nữ  trong $18$ nữ có $\mathrm{C}_{18}^1=18$ cách.\\
	Theo quy tắc nhân, số cách chọn  một đôi nam nữ là $\mathrm{C}_{20}^1\mathrm{C}_{18}^1=20\cdot 18=360$ cách.}
\end{ex}

% nguồn: MaHoaNH_DungSai_Lop10
\begin{ex}%[0D8H2-4]%[ma:e4461f93]
	Trong một lớp có $16$ bạn nữ, $18$ bạn nam. Chọn ra $5$  bạn tham gia đợt tình nguyện mùa hè xanh.
	 \choiceTF
		{\True Có $\mathrm{C}_{34}^{5}$ cách chọn $5$  bạn bất kỳ}
		{\True Có $\mathrm{C}_{16}^{5}$ cách chọn $5$  bạn đều là nữ}
		{Có $\mathrm{C}_{16}^{2}+\mathrm{C}_{18}^{3}$ cách  chọn $5$ bạn trong đó có $3$ bạn nữ, $2$ bạn nam}
		{Có $\mathrm{C}_{16}^{1}+\mathrm{C}_{18}^{4}$ cách chọn $5$ bạn  trong đó có $1$ bạn nữ, $4$ bạn nam}
	 \loigiai{
		\begin{itemchoice}
			\itemch Chọn  $5$ bạn trong $34$ bạn là $\mathrm{C}_{34}^{5}$.
			\itemch Chọn $5$ bạn nữ trong $16$  bạn là $\mathrm{C}_{16}^{5}$.
			\itemch Chọn $5$ bạn trong đó có $3$  bạn nữ, $2$ bạn nam là $\mathrm{C}_{16}^{3}\cdot \mathrm{C}_{18}^{2}$.
			\itemch Chọn $5$ bạn trong đó có $1$ bạn nữ, $4$ bạn  nam là $\mathrm{C}_{16}^{1}\cdot \mathrm{C}_{18}^{4}$.
		\end{itemchoice}
	} 
\end{ex}

% nguồn: MaHoaNH_DungSai_Lop10
\begin{ex}%[0D8H2-4]%[ma:c92642f5]
	Lớp $10$A$1$ trường THPT TenLơMan có $44$ học sinh, trong đó  có $18$ học sinh sinh nữ và $26$ học sinh nam.
	\choiceTF 
		{Số cách sắp xếp $18$ học sinh sinh nữ và $26$ học sinh  nam thành một hàng dọc là $18!\cdot 26!$}
		{Có $4!$ cách  xếp khác nhau cho $4$ học sinh ngồi vào $6$ chỗ trên một bàn dài}
		{Giáo viên chủ nhiệm cần lập ban cán sự lớp gồm $1$ lớp trưởng, $1$  lớp phó học tập và $1$ lớp phó kỉ luật. Có $3!$ cách lập}
		{\True Lớp trưởng cần chọn ra $5$ bạn và xếp thứ tự thuyết trình cho $5$ bạn này. Có $\mathrm{A}_{44}^5$ cách chọn} 
	\loigiai{ 
		\begin{itemchoice}
			 \itemch
			Có $44!$ cách sắp xếp $18$ học sinh sinh nữ và $26$ học sinh nam  thành một hàng dọc.
			 \itemch
			Có $\mathrm{A}_6^4$ cách xếp  khác nhau cho $4$ học sinh ngồi vào $6$ chỗ trên một bàn dài.
			\itemch 
			Có $\mathrm{A}_{44}^3$ cách lập  ban cán sự lớp gồm $1$ lớp trưởng, $1$ lớp phó học tập và $1$ lớp phó kỉ luật.
			\itemch Lớp trưởng cần chọn ra  $5$ bạn và xếp thứ tự thuyết trình cho $5$ bạn này. Có $\mathrm{A}_{44}^5$ cách chọn
		\end{itemchoice}
	} 
\end{ex}

% nguồn: MaHoaNH_DungSai_Lop10
\begin{ex}%[0D8H2-4]%[ma:0d39e26b]
	Từ một nhóm học sinh có $10$ nam và $15$  nữ người ta chọn ra một số học sinh.
	\choiceTF 
		{\True  Số cách chọn $2$ học sinh là $300$}
		{Số cách chọn $5$ học sinh sao cho luôn có ít nhất  $1$ học sinh nữ là $50127$}
		{Số cách chọn $5$ học sinh  sao cho luôn có ít nhất $1$ học sinh nam là $52878$}
		{\True  Số cách chọn $6$ học sinh sao cho luôn có cả nam và nữ là $171885$}
	\loigiai{ 
		\begin{itemize}
			\item Số cách chọn $2$  học sinh là $\mathrm{C}_{25}^2 = 300$.
			\item Số cách chọn $5$ học sinh sao cho luôn có ít nhất $1$ học sinh nữ là $\mathrm{C}_{25}^5 -  \mathrm{C}_{10}^5 = 52878$.
			\item Số cách chọn  $5$ học sinh sao cho luôn có ít nhất $1$ học sinh nam là $\mathrm{C}_{25}^5 - \mathrm{C}_{15}^5 = 50127$.
			\item Số cách chọn $6$ học sinh sao cho luôn có cả nam  và nữ là $\mathrm{C}_{25}^6 -\mathrm{C}_{10}^6- \mathrm{C}_{15}^6 = 171885$.
		\end{itemize}
		\begin{itemchoice}
			\itemch  Đúng
			\itemch Sai 
			\itemch  Sai
			\itemch Đúng 
		\end{itemchoice}
	 }
\end{ex}

% nguồn: MaHoaNH_DungSai_Lop10
\begin{ex}%[0D8V2-4]%[ma:cb4abde3]
	Có $3$ học sinh nam và $2$ học sinh nữ xếp  vào một bàn dài.
	\choiceTF 
		{\True Có  $120$ cách xếp các học sinh ngồi tùy ý}
		{\True Có $12$ cách xếp sao cho học sinh nam và học sinh nữ ngồi  xen kẽ}
		 {Có $24$ cách xếp sao cho học sinh nữ luôn ngồi cạnh nhau}
		{Có $24$ cách xếp sao cho học sinh nữ luôn ngồi hai đầu  bàn}
	\loigiai{ 
		Đánh số chỗ  ngồi lần lượt là $1$, $2$, $3$, $4$, $5$.
		\begin{itemchoice}
			\itemch Xếp  $5$ học sinh vào $5$ chỗ ngồi có $5!=120$ cách
			\itemch Có $3$ học sinh nam và $2$ học sinh nữ ngồi xen kẽ nhau nên học sinh nam phải ngồi  ở vị trí số $1$, $3$, $5$ và học sinh nữ ngồi ở vị trí $2$, $4$\\
			Xếp $3$ học sinh nam có $3!$ cách, xếp $2$ học sinh nữ  có $2!$ cách.\\
			Như vậy có  $3!\cdot 2! = 12$ cách.
			\itemch Xếp $2$ học sinh nữ có  $2!$ cách.\\
			Coi $2$ học sinh nữ là  một phần tử $x$, xếp $x$ và $3$ học sinh nam có $4!$ cách.\\
			Như vậy có $2!\cdot  4!=48$ cách.
			\itemch Xếp $2$ học sinh nữ vào vị trí  số $1$ và số $5$ có $2!$ cách\\
			Xếp $3$ học sinh vào các vị trí $2$, $3$, $4$ có  $3$! cách.\\
			Như  vậy có $2!\cdot3!=12$ cách.
		\end{itemchoice}
	 }
\end{ex}


%%%%%%%%%%%%%%%% TRẢ LỜI NGẮN %%%%%%%%%%%%%%%%

% nguồn: MaHoaNH_TraLoiNgan_Lop10
\begin{bt}%[0D8H2-4]%[ma:56ae816f]
	Cô giáo chọn trong lớp $3$ bạn học sinh là An, Bình và Cường để làm ban cán sự lớp gồm lớp trưởng, lớp phó học tập, lớp phó lao động. Hỏi cô có bao nhiêu cách phân  công $3$ bạn làm các chức vụ trên.
	 \shortans{6}
	 \loigiai{
		Mỗi cách phân công $3$ bạn An, Bình và Cường để làm các chức vụ lớp trưởng, lớp phó học tập, lớp phó lao động  là một hoán vị của $3$ bạn trên.\\
		Do đó  số cách phân công là $\mathrm{P}_{3}=3!=6$.
	 }
\end{bt}

% nguồn: MaHoaNH_TraLoiNgan_Lop10
\begin{bt}%[0D8H2-4]%[ma:a7b467d9]
	 Cần xếp một nhóm $4$ học sinh ngồi vào một dãy $4$ chiếc ghế. Nếu bạn Nga (một thành viên trong nhóm) nhất định muốn ngồi vào chiếc ghế ngoài cùng bên trái, thì có bao nhiêu cách xếp?
	\shortans{6} 
	\loigiai{ 
		Ta có Nếu bạn Nga ngồi vào ghế ngoài cùng bên trái, thì cần xếp ba bạn còn lại vào $3$ ghế còn lại, số cách sắp xếp chính là số các  hoán vị của $3$ phần tử, nên có $3!$ cách xếp.\\
		Vậy có $6$ cách sắp xếp  $4$ bạn vào bốn ghế ngồi với ghế ngoài cùng bên trái bạn Nga ngồi.
	 }
\end{bt}

% nguồn: MaHoaNH_TraLoiNgan_Lop10
\begin{bt}%[0D8H2-4]%[ma:11231bf4]
	Đội văn nghệ của nhà trường gồm $4$ học sinh lớp 12A, $3$ học sinh lớp 12B và $2$ học sinh lớp $12$. Chọn ngẫu nhiên $5$ học sinh từ đội văn nghệ để biễu diễn trong lễ bế giảng. Hỏi có bao nhiêu cách chọn sao cho lớp nào cũng có học sinh được  chọn?
	\shortans{98} 
	\loigiai{ 
		Số cách chọn  ngẫu nhiên $5$ học sinh $\mathrm{C}_{9}^{5}$ cách.\\
		Số  cách chọn $5$ học sinh chỉ có $2$ lớp $\mathrm{C}_{7}^{5}+\mathrm{C}_{6}^{5}+\mathrm{C}_{5}^{5}$\\
		 Vậy số cách chọn $5$ học sinh có cả $3$ lớp là $\mathrm{C}_{9}^{5}-\left(\mathrm{C}_{7}^{5}+\mathrm{C}_{6}^{5}+\mathrm{C}_{5}^{5}\right)=98$.
	 }
\end{bt}

% nguồn: MaHoaNH_TraLoiNgan_Lop10
\begin{bt}%[0D8H2-4]%[ma:8d3f9c2b]
	Có bao nhiêu  cách sắp xếp $5$ bạn nam và $5$ bạn nữ ngồi xung quanh một bàn tròn sao cho nam và nữ ngồi xen kẽ?
	 \shortans{2880}
	\loigiai{ 
		Chọn $1$ bạn nam ngồi cố định vào $1$ vị trí, $9$ bạn còn lại sẽ hoán vị xung quanh bạn này theo nguyên tắc là nam nữ xen kẽ. Khi đó, các bạn nam còn lại sẽ ở  vị trí mang số $3$, $5$, $7$, $9$ và nữ sẽ ở vị trí số $2$, $4$, $6$, $8$, $10$ (theo chiều kim đồng hồ).\\
		Ở mỗi vị trí của  mình, các nam và nữ được hoán vị cho nhau.\\
		Do đó, có $1\cdot 4!\cdot 5!=2880$ cách  sắp xếp.
	} 
\end{bt}

% nguồn: MaHoaNH_TraLoiNgan_Lop10
\begin{bt}%[0D8H2-4]%[ma:a345c86e]
	Một nhóm học sinh  gồm $4$ học sinh nam và $5$ học sinh nữ. Hỏi có bao nhiêu cách sắp xếp $9$ học sinh trên thành $1$ hàng dọc sao cho nam nữ đứng xen kẽ?
	 \shortans{2880}
	 \loigiai{
		Xếp $4$ học sinh nam thành  hàng dọc có $4!$ cách xếp.\\
		Giữa $4$ học sinh nam có $5$ khoảng trống ta xếp các bạn nữ vào  vị trí đó nên có $5!$ cách xếp.\\
		 Theo quy tắc nhân có $4!\cdot 5!=2880$ cách xếp thoả mãn bài ra.
	 }
\end{bt}

% nguồn: MaHoaNH_TraLoiNgan_Lop10
\begin{bt}%[0D8H2-4]%[ma:fd5844de]
	Một đội thanh niên tình  nguyện có $9$ người gồm $6$ nam và $3$ nữ.
	Hỏi có bao nhiêu cách phân công đội thanh niên tình nguyện đó về $3$ tỉnh miền núi sao  cho mỗi tỉnh có $3$ nam và $1$ nữ.
	\shortans{540} 
	\loigiai{Chọn $2$ nam và $1$  nữ về tỉnh thứ nhất có $\mathrm{C}_{6}^2\cdot \mathrm{C}_3^1=45$ (cách).\\
		Khi đó còn lại $4$ nam và $2$ nữ. Chọn $2$ nam và $1$ nữ về tỉnh  thứ hai có $\mathrm{C}_{4}^2\cdot \mathrm{C}_2^1=12$ (cách). \\
		Cuối cùng còn $2$ nam và $1$ nữ nên có $1$ cách chọn $2$ nam và $1$  nữ này về tỉnh thứ ba.\\
		Vậy số cách phân công thỏa mãn yêu cầu  bài toán là
	$45\cdot 12\cdot  1=540$.}
\end{bt}

% nguồn: MaHoaNH_TraLoiNgan_Lop10
\begin{bt}%[0D8H2-4]%[ma:11497bf7]
	Một nhóm học sinh gồm $10$ học sinh nam và $7$ học sinh nữ. Từ nhóm này, hỏi có bao nhiêu cách chọn ra một nhóm nhỏ gồm $4$  học sinh sao cho nhóm có đủ cả nam và nữ?
	\shortans{2135} 
	\loigiai{ 
		\begin{itemize}
			\item Chọn $4$ học sinh bất kì có  $\mathrm{C}_{17}^4$ cách.
			\item Chọn $4$ học sinh nam có $\mathrm{C}_{10}^4$ cách. 
			\item Chọn $4$ học  sinh nữ có $\mathrm{C}_7^4$ cách.
		\end{itemize}
		Vậy  số cách chọn $4$ học sinh có ít nhất $1$ học sinh nữ là $\mathrm{C}_{17}^4 - \mathrm{C}_{10}^4 - \mathrm{C}_7^4 = 2\,135$.
	} 
\end{bt}

% nguồn: MaHoaNH_TraLoiNgan_Lop10
\begin{bt}%[0D8H2-4]%[ma:5dadfa08]
	Đội thanh niên xung kích của một trường phổ thông có $12$ học sinh, gồm $5$ học sinh khối $10$, $4$ học sinh khối $11$ và $3$ học sinh khối $12$. Cần chọn $4$ học sinh đi làm nhiệm vụ, sao cho $4$ học sinh này thuộc không quá $2$ trong $3$ khối trên. Hỏi  có bao nhiêu cách chọn như vậy?
	 \shortans{225}
	\loigiai{ 
		Số cách chọn $4$ học sinh mà mỗi khối có ít nhất một em được tính như  sau
		\begin{itemize}
			\item Khối $10$ có $2$ học  sinh, các khối $11$, $12$ mỗi khối có $1$ học sinh. \\
			Số cách chọn  là $\mathrm{C}_5^2\cdot \mathrm{C}_4^1\cdot \mathrm{C}_3^1=120$.
			\item  Khối $11$ có $2$ học sinh, các khối $10$, $12$ mỗi khối có $1$ học sinh. \\
			Số cách chọn là $\mathrm{C}_5^1\cdot \mathrm{C}_4^2\cdot  \mathrm{C}_3^1=90$.
			\item Khối $12$ có $2$ học sinh, các khối $10$, $11$ mỗi khối  có $1$ học sinh. \\
			Số cách chọn  là $\mathrm{C}_5^1\cdot \mathrm{C}_4^1\cdot \mathrm{C}_3^2=60$.
		\end{itemize}
		Suy ra, số cách chọn $4$ học sinh mà mỗi  lớp có ít nhất một học sinh là $120+90+60=270 $.\\
		Vậy số cách chọn phải tìm là $495-270=225 $. 
	 }
\end{bt}

% nguồn: MaHoaNH_TraLoiNgan_Lop10
\begin{bt}%[0D8V2-4]%[ma:ff7ee978]
	Trong một tổ học sinh có $6$ học sinh nữ $10$ học sinh nam. Hạnh là một  trong $6$ học sinh nữ, Huy là một trong $10$ học sinh nam. Cô chủ nhiệm cần chọn ra $5$ bạn trong tổ để tham gia hoạt động văn nghệ nhân ngày 20/11 sắp tới. Hỏi cô chủ nhiệm có bao nhiêu cách chọn trong đó có ít nhất $2$ em Hạnh và Huy không được tham gia.
	\shortans{4004} 
	\loigiai{ 
		Bài toán: dùng biến cố đối tìm số cách chọn ra $5$ bạn  trong đó có Hạnh và Huy.
		\begin{itemize}
			\item \textbf{Bước 1:} Chọn nhóm $3$ học sinh  trong $14$ học sinh (không tính Huy và Hạnh): $\mathrm{C}_{14}^{3}=364$ cách.
			\item \textbf{Bước 2:} Chọn  $2$ học sinh Huy và Hạnh có $1$ cách.
		\end{itemize}
		Do đó, có $364$ cách chọn  ra $5$ học sinh trong tổ mà có hai học sinh Huy và Hạnh.\\
		Chọn ra $5$ học sinh bất kì trong $16$  học sinh $\mathrm{C}_{16}^{5}=4368$ cách.\\
		 Vậy cô chủ nhiệm có số cách chọn trong đó có ít nhất $2$ em Hạnh và Huy không được tham gia là
		 $4368-364=4004$ (cách).
	} 
\end{bt}

% nguồn: MaHoaNH_TraLoiNgan_Lop10
\begin{bt}%[0D8V2-4]%[ma:3678c5ba]
	Ở một  Đoàn trường phổ thông có $5$ thầy giáo, $4$ cô giáo và $8$ học sinh. Có $X$ cách chọn ra một đoàn công tác gồm $7$ người trong đó có $1$ trưởng đoàn là thầy giáo, $1$ phó đoàn là cô giáo và đoàn công tác có ít nhất $4$ học sinh. Tính tổng các chữ số của $X$.
	 \shortans{12}
	\loigiai{ 
		\begin{itemize}
			\item Trường hợp $1$: Đoàn công tác gồm $1$ thầy  giáo, $1$ cô giáo và $5$ học sinh.
			\begin{itemize}
				\item  Chọn ra $1$ thầy giáo từ $5$ thầy giáo có $\mathrm{C}_5^1$ cách.
				\item  Chọn ra $1$ cô giáo từ $4$ cô giáo có $\mathrm{C}_4^1$ cách.
				\item Chọn ra $5$  học sinh từ $8$ học sinh có $\mathrm{C}_8^5$ cách.
			\end{itemize}
			Do chỉ có $1$ thầy giáo và $1$ cô giáo nên chỉ có $1$ cách chọn  trưởng đoàn và phó đoàn.\\
			Theo  quy tắc nhân có $\mathrm{C}_5^1\cdot\mathrm{C}_4^1\cdot\mathrm{C}_8^5\cdot1=1 \, 120$ cách.
			\item Trường hợp $2$: Đoàn công tác gồm $2$ thầy  giáo, $1$ cô giáo và $4$ học sinh.
			\begin{itemize}
				\item Chọn ra $2$ thầy giáo từ $5$ thầy  giáo có $\mathrm{C}_5^2$ cách.
				\item Chọn ra $1$ cô giáo từ $4$  cô giáo có $\mathrm{C}_4^1$ cách.
				\item  Chọn ra $4$ học sinh từ $8$ học sinh có $\mathrm{C}_8^4$ cách.
				\item Chọn $1$ thầy giáo từ $2$ thầy giáo  làm trưởng đoàn có $\mathrm{C}_2^1$ cách.
			\end{itemize}
			Theo quy tắc nhân có $\mathrm{C}_5^2\cdot\mathrm{C}_4^1\cdot\mathrm{C}_8^4\cdot\mathrm{C}_2^1=5\,600$  cách.
			\item Trường hợp $3$: Đoàn  công tác gồm $1$ thầy giáo, $2$ cô giáo và $4$ học sinh.
			\begin{itemize}
				 \item Chọn ra $1$ thầy giáo từ $5$ thầy giáo có $\mathrm{C}_5^1$ cách.
				\item Chọn ra $2$ cô giáo từ $4$ cô giáo  có $\mathrm{C}_4^2$ cách.
				\item Chọn ra $4$ học sinh từ $8$ học sinh  có $\mathrm{C}_8^4$ cách.
				\item  Chọn $1$ cô giáo từ $2$ cô giáo làm phó đoàn có $\mathrm{C}_2^1$.
			\end{itemize}
			Theo quy tắc nhân có $\mathrm{C}_5^1\cdot\mathrm{C}_4^2\cdot\mathrm{C}_8^4\cdot\mathrm{C}_2^1=4\, 200$  cách.
		\end{itemize}
		Theo quy tắc cộng ta có $1\,120+5\,600+4\,200=10\,920$  cách.
	} 
\end{bt}


%%%%%%%%%%%%%%%% TỰ LUẬN %%%%%%%%%%%%%%%%

% nguồn: MaHoaNH_TuLuan_Lop10
\begin{bt}%[0D8H2-4]%[ma:196db9f3]
	Một lớp học có ba cán bộ lớp là $A$, $B$, $C$. Có bao  nhiêu cách chọn một lớp trưởng, một lớp phó, một bí thư từ ba cán bộ lớp $A$, $B$, $C$?
	 \loigiai{
		Mỗi một cách chọn là  một hoán vị của $3$ phần tử.\\
		Do đó có  $\mathrm{P}_3 = 3! = 6$ cách chọn.
	} 
\end{bt}

% nguồn: MaHoaNH_TuLuan_Lop10
\begin{bt}%[0D8H2-4]%[ma:8a610984]
	Một lớp có $20$ học sinh nam, $15$ học sinh nữ. Có bao nhiêu cách để giáo viên chủ nhiệm chọn ra một ban chấp hành Đoàn $3$ người  sao cho ban chấp hành có ít nhất $1$ nữ.
	 \loigiai{
		Số  cách chọn $3$ học sinh tùy ý từ $35$ học sinh là: $N = \mathrm{C}_{35}^3 = 6545$ cách.\\
		Số cách chọn $3$ học sinh không có nữ là: $\overline{n} = \mathrm{C}_{20}^3 =  1140$ cách.\\
		Vậy số cách chọn $3$ học sinh  có ít nhất $1$ nữ là: $n = N - \overline{n} = \mathrm{C}_{35}^3 - \mathrm{C}_{20}^3 = 5405$ cách.
	 }
\end{bt}

% nguồn: MaHoaNH_TuLuan_Lop10
\begin{bt}%[0D8H2-4]%[ma:466e9e01]
	Một đội cảnh sát giao thông gồm $15$ người trong đó có $12$ nam. Hỏi có bao nhiêu  cách phân đội cảnh sát giao thông đó về $3$ chốt giao thông sao cho mỗi chốt có $4$ nam và $1$ nữ.
	\loigiai{ 
		Chọn $4$ nam trong $12$ nam và chọn $1$  nữ trong $3$ nữ, có $\mathrm{C}_{12}^4\cdot \mathrm{C}_3^1$ cách. \\
		Chọn $4$ nam trong $8$ nam còn lại và chọn $1$ nữ trong $2$ nữ còn lại, có  $\mathrm{C}_8^4\cdot \mathrm{C}_2^1$ cách. \\
		Chọn $4$ nam còn lại và $1$ nữ còn lại bắt buộc phải về công tác ở chốt giao  thông cuối cùng, nên có $1$ cách. \\
		Theo quy tắc  nhân có $\mathrm{C}_{12}^4\cdot \mathrm{C}_3^1\cdot \mathrm{C}_8^4\cdot \mathrm{C}_2^1\cdot 1=207\,900$ cách chọn.
	 }
\end{bt}

% nguồn: MaHoaNH_TuLuan_Lop10
\begin{bt}%[0D8H2-4]%[ma:468a2e66]
	Một trường trung học phổ thông có $8$ giáo viên Toán gồm có $3$ nữ và $5$ nam, giáo viên Vật lý thì có $4$  giáo viên nam (không có giáo viên nữ). Nhà trường muốn chọn ra một đoàn thanh tra công tác ôn thi THPTQG gồm $3$ giáo viên sao cho phải có cả giáo viên Toán và Vật lý, đồng thời phải có cả nam và nữ. Hỏi nhà trường có bao nhiêu cách chọn ra một ban thanh tra?
	 \loigiai{
		Vì chọn ra $3$ người mà yêu  cầu phải có giáo viên nam và giáo viên nữ trong đoàn nên số giáo viên nữ được chọn chỉ có thể bằng $1$ hoặc $2$. Ta xét hai trường hợp sau.\\
		{\bfseries Trường hợp 1.} Chọn $1$ giáo viên nữ:  có $\mathrm{C}_3^1$ cách. Khi đó,\\
		Chọn $1$ giáo viên nam môn Toán  và $1$ nam môn Vật lý: có $\mathrm{C}_5^1\cdot \mathrm{C}_4^1$ cách.\\
		Chọn $2$  giáo viên nam môn Vật lý: có $\mathrm{C}_4^2$ cách.\\
		Trường hợp này có $\mathrm{C}_3^1\left(\mathrm{C}_5^1\cdot \mathrm{C}_4^1+\mathrm{C}_4^2\right)$ cách chọn.\\ 
		{\bfseries Trường hợp 2.} Chọn $2$ giáo viên nữ: có $\mathrm{C}_3^2$ cách chọn. Khi đó, chọn thêm $1$ giáo  viên nam môn Vật lý: có $\mathrm{C}_4^1$ cách.\\
		Trường hợp  này có $\mathrm{C}_3^2\cdot \mathrm{C}_4^1$ cách chọn.\\
		Vậy tất cả có $\mathrm{C}_3^1\left(\mathrm{C}_5^1\cdot  \mathrm{C}_4^1+\mathrm{C}_4^2\right)+\mathrm{C}_3^2\cdot \mathrm{C}_4^1=90$ cách chọn.
	 }
\end{bt}

% nguồn: MaHoaNH_TuLuan_Lop10
\begin{bt}%[0D8H2-4]%[ma:f24aab80]
	Một đội văn nghệ gồm $20$ người, trong đó có $10$  nam, $10$ nữ. Hỏi có bao nhiêu cách chọn ra $5$ người, sao cho
	\begin{enumerate}
		\item Có đúng $2$ nam trong $5$  người đó.
		\item Có ít nhất $2$ nam, ít nhất  $1$ nữ trong $5$ người đó.
	\end{enumerate}
	 \loigiai{
		\begin{enumerate}
			\item Số cách chọn $2$ nam, $3$ nữ là $\mathrm{C}_{10}^2\mathrm{C}_{10}^3=5\,400$  cách.
			\item Có các  trường hợp xảy ra thỏa yêu cầu của đề như sau
			\begin{itemize}
				\item Trường hợp 1: Có $2$ nam và $3$ nữ. Số cách  chọn $5\,400$ cách.
				\item Trường hợp 2: Có $3$ nam và $2$ nữ. Số cách chọn $\mathrm{C}_{10}^3\mathrm{C}_{10}^2=5\,400$  cách.
				\item Trường hợp 3: Có $4$ nam và $1$ nữ.  Số cách chọn $\mathrm{C}_{10}^4\mathrm{C}_{10}^1=2\,100$ cách.
			\end{itemize}
			Tổng cộng $3$ trường hợp ta  có $5\,400+5\,400+2\,100=12\,900$ cách.
		\end{enumerate}
	 }
\end{bt}

% nguồn: MaHoaNH_TuLuan_Lop10
\begin{bt}%[0D8H2-4]%[ma:fb0fd7a1]
	Môt lớp có $20$ học sinh trong đó có $14$ nam, $6$ nữ. Hỏi có bao nhiêu cách lập $1$ đội  gồm $4$ học sinh trong đó có
	\begin{enumerate}
		\item Số nam  và nữ bằng nhau.
		\item Có  ít nhất 1 nữ.
	\end{enumerate}
	 \loigiai{
		\begin{enumerate}
			\item Chọn $2$ nam  trong $14$ nam, có $\mathrm{C}_{14}^2$ cách. \\
			Chọn $2$ nữ trong  $6$ nữ, có $\mathrm{C}_6^2$ cách. \\
			Vậy số cách chọn nhóm có $2$ nam,  $2$ nữ là $\mathrm{C}_{14}^2\cdot \mathrm{C}_6^2=1\,365$ cách.
			\item \textbf{Cách 1:} Xét các  trường hợp xảy ra cụ thể:
			\begin{itemize}
				\item Trường hợp 1: Chọn $1$ nữ, $3$  nam có $6\cdot \mathrm{C}_{14}^3=2\,184$ cách.
				\item Trường hợp 2: Chọn  $2$ nữ, $2$ nam có $\mathrm{C}_{14}^2\cdot \mathrm{C}_6^2=1365$ cách.
				\item Trường hợp 3: Chọn $3$ nữ, $1$ nam  có $\mathrm{C}_6^3\cdot 14=280$ cách.
				 \item Trường hợp 4: Chọn $4$ nữ thì có $\mathrm{C}_6^4=15$ cách.
			\end{itemize}
			Vậy số cách chọn cần tìm là $2184+1365+280+15=3\,844$  cách. \\
			 \textbf{Cách 2:} Sử dụng phần bù \\
			Chọn $4$ bạn bất kỳ trong $20$ bạn, có  $\mathrm{C}_{20}^4$ cách. \\
			Chọn $4$  bạn đều nam, có $\mathrm{C}_{14}^4$ cách. \\
			Suy ra chọn $4$  bạn có ít nhất $1$ nữ: $\mathrm{C}_{20}^4-\mathrm{C}_{14}^4=3\,844$ cách chọn.
		\end{enumerate}
	} 
\end{bt}
